\section{Host software}
\label{sec:host}

The host side (\file{sw/gemm\_host/}, C++17) is compiled per build against that build's generated \code{csr.h} and
\code{liblitepcie}, so register addresses always match the bitstream. Three small classes carry everything:
\begin{description}[font=\sffamily\small\bfseries, leftmargin=1.2em]
  \item[\code{Card}] opens \code{/dev/litepcie0}, reads and writes CSRs through the driver, maps both DMA rings
        (zero-copy) and enables or disables the DMA engines.
  \item[\code{DdrDma}] moves data between host memory and a DDR channel through the bridge in 2\,MB chunks, with the
        protocol of \cref{sec:pcie:rates}, and reports where the time went (copy, start, wait, stop).
  \item[\code{GemmKernel}] writes the kernel arguments, starts it, polls \code{done}, and returns the cycle count, the
        measured accel clock and every performance counter.
\end{description}

\code{gemm\_runner} is the benchmark program built on them; \code{ddr\_dma\_test} checks and times the transfers on
their own. A run of \code{gemm\_runner} (\cref{fig:sequence}):
\begin{enumerate}
  \item checks the bitstream's \code{build\_id} against the expected configuration hash, the matrix sizes against the
        build's size rules (\code{--multiples}: whole 64-byte words, e.g.\ \(K\) and \(M\) multiples of 16 floats,
        and \(N\) too with transposed A) and the matrix placement against the channel capacity;
  \item generates \(A\) and \(B\) from a seed and uploads \(A^\mathsf{T}\) (or \(A\)) and \(B\) to their channels;
  \item for each warm-up and measured repetition, runs the kernel and downloads \(C\);
  \item verifies \(C\) against a double-precision host reference, either every element (\code{full}) or 4096 random
        elements plus two corners (\code{sampled}), with the tolerance
        \(|c - c_\mathrm{ref}| \le 4K\varepsilon\sum_k|a_{ik}b_{kj}|\), \(\varepsilon = 2^{-23}\);
  \item prints one JSON row per repetition: sizes, cycles, measured clock, kernel time and GFLOP/s, transfer bytes and
        times, end-to-end time, verification result and all counters.
\end{enumerate}
Kernel time is \(\text{cycles}/f_\text{accel}\), with the clock measured on the card, so it does not include the
host's polling. Sizes the kernel cannot handle used to give wrong results without any error (e.g.\ \(M\) not a
multiple of 16 floats); the size check now rejects them before any transfer.

\begin{figure}[t]
  \centering
  \resizebox{\linewidth}{!}{% One GEMM run as seen by gemm_runner: build-id check, uploads, kernel, download, verification.
\begin{tikzpicture}[x=1cm,y=1cm, every node/.append style={font=\sffamily\footnotesize}]
  \tikzset{life/.style={draw=black!35, densely dashed},
           msg/.style={->, thick, draw=black!75},
           act/.style={fill=#1, draw=black!50, minimum width=2.2mm, inner sep=0pt},
           m/.style={lbl, fill=white, fill opacity=0.9, text opacity=1, inner sep=1pt}}
  \def\top{0} \def\bot{-11.2}
  \foreach \x/\n/\s in {0/gemm\_runner/hst, 3.6/litepcie.ko/hst, 7.4/DMA + bridge/sys, 11.0/GEMM kernel/acc, 14.4/{DDR3 CH0 + CH1}/ddr} {
    \node[\s, minimum width=25mm] at (\x,\top) {\n};
    \draw[life] (\x,\top-0.35) -- (\x,\bot);
  }
  % 1. check
  \draw[msg] (0,-0.9) -- node[m, above] {ioctl REG: read \code{build\_id}} (3.6,-0.9);
  \draw[msg] (3.6,-1.15) -- node[m, above] {MMIO} (7.4,-1.15);
  \node[note, anchor=west, text width=30mm] at (14.9,-1.0) {refuses to run if it is not the build's config hash};
  % 2. uploads
  \fill[cBlue!8] (-1.3,-1.6) rectangle (16.9,-4.4);
  \node[lbl, anchor=north east, font=\sffamily\scriptsize\bfseries] at (16.9,-1.6) {upload A\(^\mathsf{T}\) (CH0) and B (CH1), per 2\,MB chunk};
  \draw[msg] (0,-2.2) -- node[m, above] {copy chunk into the H2C ring (mmap)} (3.6,-2.2);
  \draw[msg] (0,-2.65) -- node[m, above, pos=0.73] {bridge: sel, writer base/length, enable} (7.4,-2.65);
  \draw[msg] (0,-3.1) -- node[m, above] {ioctl: DMA reader enable} (3.6,-3.1);
  \draw[msg, cBlue!70!black] (7.4,-3.35) -- node[m, above, pos=0.45] {MRd \(\rightarrow\) completions} (3.6,-3.35);
  \draw[msg, cBlue!70!black] (7.4,-3.55) -- node[m, below] {512\,b writes} (14.4,-3.55);
  \draw[msg] (0,-4.0) -- node[m, above, pos=0.73] {poll \code{drained}; reader disable (driver waits 1\,ms)} (7.4,-4.0);
  % 3. kernel
  \fill[cOrange!10] (-1.3,-4.6) rectangle (16.9,-7.2);
  \node[lbl, anchor=north east, font=\sffamily\scriptsize\bfseries] at (16.9,-4.6) {kernel};
  \draw[msg] (0,-5.15) -- node[m, above, pos=0.8] {CSRs: a, b, c, size\_n/k/m; \code{start}} (11.0,-5.15);
  \node[act=cOrange!40, minimum height=16mm, anchor=north] at (11.0,-5.25) {};
  \draw[msg, cOrange!70!black] (11.0,-5.6) -- node[m, above] {AXI reads A, B} (14.4,-5.6);
  \draw[msg, cOrange!70!black] (14.4,-5.95) -- (11.0,-5.95);
  \draw[msg, cOrange!70!black] (11.0,-6.35) -- node[m, above] {AXI writes C} (14.4,-6.35);
  \draw[msg] (0,-5.7) -- node[m, above, pos=0.35] {poll \code{status.done}} (11.0,-5.7);
  \draw[msg] (11.0,-6.85) -- node[m, above, pos=0.2] {\code{cycles}, perf counters (after done)} (0,-6.85);
  % 4. download
  \fill[cGreen!8] (-1.3,-7.4) rectangle (16.9,-10.0);
  \node[lbl, anchor=north east, font=\sffamily\scriptsize\bfseries] at (16.9,-7.4) {download C (CH1): writer enabled once, then per chunk};
  \draw[msg] (0,-7.95) -- node[m, above] {ioctl: DMA writer enable} (3.6,-7.95);
  \draw[msg] (0,-8.4) -- node[m, above, pos=0.73] {bridge: reader base/length, enable} (7.4,-8.4);
  \draw[msg, cGreen!60!black] (14.4,-8.75) -- node[m, above] {512\,b reads} (7.4,-8.75);
  \draw[msg, cGreen!60!black] (7.4,-9.0) -- node[m, below, pos=0.45] {MWr into C2H ring} (3.6,-9.0);
  \draw[msg] (0,-9.6) -- node[m, above, pos=0.73] {wait: writer loop status \(\ge\) buffers; copy out} (7.4,-9.6);
  % 5. verify
  \draw[msg] (0,-10.4) -- ++(0.6,0) -- ++(0,-0.35) -- ++(-0.55,0);
  \node[note, anchor=west, text width=60mm] at (0.75,-10.55) {verify against the host reference (full or sampled), print one JSON row per repetition};
\end{tikzpicture}
}
  \caption{One GEMM run: build-id check, uploads in 2\,MB chunks, the kernel, the download of C, verification.}
  \label{fig:sequence}
\end{figure}
